\section{预期成果与创新点}\label{sec:outcomes}

\subsection{预期成果概述}

\begin{tcolorbox}[outcomestyle,title=成果一：开源ECC AI全自动设计框架E3A]
发布第一个端到端纠错码AI全自动设计开源框架，涵盖从需求解析到RTL代码输出的全部模块。框架支持LDPC（BG1/BG2）与Polar（CA-Polar、PAC码）两类主流码型，信道模型覆盖AWGN、瑞利衰落、Flash双读取、量子Pauli信道与忆阻器非参数信道五大类。预期在GitHub获得显著学术关注度，推动该领域开放性研究。
\end{tcolorbox}

\begin{tcolorbox}[outcomestyle,title=成果二：高精度ECC性能预测神经算子模型]
发布预训练DeepONet性能预测模型库（覆盖五大信道类型），实现：
\begin{itemize}[nosep,leftmargin=1.5em]
  \item BER预测误差 $<0.1$~dB（相较蒙特卡洛$\mathrm{BER}=10^{-6}$基准）
  \item FER预测误差 $<0.15$~dB
  \item 单次推理时间 $<1$~ms（GPU）
  \item 相较蒙特卡洛仿真加速比 $>10^3\times$
\end{itemize}
\end{tcolorbox}

\begin{tcolorbox}[outcomestyle,title=成果三：跨工艺VLSI性能代理模型库]
发布涵盖28~nm与7~nm工艺节点的LDPC/Polar译码器硬件代理模型，实现面积预测误差$<5\%$、功耗预测误差$<8\%$、时延预测误差$<6\%$（相较Design Compiler综合结果）。训练数据集（$>5000$组综合结果）与模型权重同步开源。
\end{tcolorbox}

\begin{tcolorbox}[outcomestyle,title=成果四：高水平学术发表]
预期发表：
\begin{itemize}[nosep,leftmargin=1.5em]
  \item 顶级期刊论文（\textit{Nature}/$\textit{Nature Electronics}$/$\textit{Nature Communications}$）1篇，聚焦神经算子ECC性能预测的理论与实验
  \item 顶级会议论文（\textit{ISSCC}/\textit{DAC}/\textit{ICCAD}）2--3篇，覆盖硬件代理模型、LLM需求解析与联合优化
  \item 专利申请（ECC AI全自动设计方法）1项
\end{itemize}
\end{tcolorbox}

\subsection{量化成果指标汇总}

\begin{table}[htbp]
\centering
\caption{E3A预期量化成果指标}
\label{tab:outcomes_quant}
\setlength{\tabcolsep}{5pt}
\rowcolors{2}{lightblue!15}{white}
\begin{tabular}{@{}L{3.2cm}L{3.5cm}L{3.5cm}L{2.5cm}@{}}
\toprule
\textbf{指标} & \textbf{当前SOTA} & \textbf{E3A预期} & \textbf{提升倍率} \\
\midrule
BER预测误差 & $>0.5$~dB（启发式） & $<0.1$~dB & $>5\times$ \\
FER预测误差 & $>0.8$~dB（启发式） & $<0.15$~dB & $>5\times$ \\
性能评估速度 & 小时级（MC仿真） & $<1$~ms & $>10^6\times$ \\
设计全流程时间 & 数月（人工） & $<1$小时 & $>10^3\times$ \\
面积预测误差 & 无代理（DC综合） & $<5\%$ & \textit{首创} \\
可覆盖信道类型 & 1--2类 & $\geq 5$类 & $>2.5\times$ \\
自动化程度 & 局部辅助 & 全流程闭环 & \textit{首创} \\
RTL代码生成 & 无 & 可综合RTL & \textit{首创} \\
开放数据集 & 无 & $>5000$组 & \textit{首创} \\
\bottomrule
\end{tabular}
\end{table}

\subsection{核心创新点}

\textbf{创新点一：通用信道三元组抽象}——首次将任意物理信道统一表示为$(\{\SNR_i\},\bvec{R},p(x))$三元组，打破信道类型壁垒，使单一框架适用于五类截然不同的物理系统。

\textbf{创新点二：ECC性能神经算子}——首次将DeepONet范式引入纠错码性能预测，以严格的函数空间算子逼近理论为支撑，突破传统点估计局限，实现全SNR曲线高精度预测。

\textbf{创新点三：算法-硬件联合代理优化}——首次将算法性能代理与硬件性能代理在统一优化框架中联合使用，消除算法-硬件鸿沟，实现帕累托最优的协同设计。

\textbf{创新点四：LLM驱动的形式化需求解析}——首次在ECC设计中引入RAG+CoT的形式化需求解析链，将自然语言工程规范自动转化为香农极限感知的形式化约束集合，实现非专家可用的交互式设计入口。